\documentclass[11pt]{article}
\usepackage[utf8]{inputenc}
\usepackage{amsmath,amssymb,amsthm}
\usepackage[margin=1in]{geometry}
\usepackage{hyperref}
\usepackage{xcolor}

\newcommand{\tideref}[2][]{\href{https://therisensea.org/tides/#2/}{\texttt{#2}}}

\title{Tide retrospective: recovery in dimension one}
\author{Tide loop (Claude Fable 5, with GPT-5.6 Sol deliberation)}
\date{2026-08-09}

\begin{document}
\maketitle

\section{Setting}

The germbij Theorem~7.3 arc established \emph{identifiability}: the
expansion data determines the loss germ. The note's Section~7.4 goes
further and \emph{recovers} the germ constructively; in dimension one,
the leading exponent of the partition function recovers the degree $2k$
of an isolated minimum and the moments recover the coefficients. This
tide formalises the smallest real instance of that claim, opening a
constructive-recovery thread on the seabed whose closed
forms\footnote{\texttt{partitionFunction\_kthPotential} and the
\texttt{kth\_moment\_even} family.} were laid down in the May tides.

\section{The thing formalised}

Three declarations in \texttt{Laplace/OneD/Recovery.lean}, sorry-free
with zero warnings. The rigidity lemma\footnote{\texttt{eventual\_power\_eq}.}:
if $\alpha_1 t^{\beta_1} = \alpha_2 t^{\beta_2}$ for all $t \geq T$ with
$\alpha_1 > 0$, then $\beta_1 = \beta_2$ and $\alpha_1 = \alpha_2$
(positivity of $\alpha_2$ is not needed). The scaling
identity\footnote{\texttt{partitionFunction\_smul}.}
$Z_{aL}(t) = Z_L(at)$. And the recovery
theorem\footnote{\texttt{kth\_partitionFunction\_recovery}.}: for
potentials $a_i\, x^{2k_i}/(2k_i)!$ with $a_i > 0$, $k_i \geq 1$, if the
partition functions agree on a ray $[T, \infty)$, $T > 0$, then
$k_1 = k_2$ and $a_1 = a_2$. The exponent recovers the degree, the
coefficient recovers the scale: the note's Section~7.4 statement for the
zeroth moment.

\section{Proof strategy}

The closed form gives $Z(t) = \alpha\, t^{-1/(2k)}$ with
$\alpha = \frac{1}{k}\bigl((2k)!/a\bigr)^{1/(2k)}\Gamma(\tfrac{1}{2k})$,
keeping $\bigl((2k)!/a\bigr)^{1/(2k)}$ as a single base per the consult's
advice (splitting it multiplies the \texttt{rpow} bookkeeping for no
benefit). For the rigidity lemma, sampling at $t_0 = \max(T, 2)$ and
$t_0 \cdot t_0$ and writing $X = t_0^{\beta_1}$, $Y = t_0^{\beta_2}$
turns the two equations into $\alpha_1 X = \alpha_2 Y$ and
$\alpha_1 X^2 = \alpha_2 Y^2$, whence
$(\alpha_1 X)(X - Y) = 0$ by a one-line
\texttt{linear\_combination}; $X = Y$ forces $\beta_1 = \beta_2$ by
strict monotonicity of the exponent at base $t_0 > 1$, and the
coefficients cancel. No logarithms anywhere. The recovery theorem is
then two injectivity steps: $k \mapsto -1/(2k)$ on positive naturals,
and the base of a positive-exponent \texttt{rpow} on nonnegative reals.

\section{Roadblocks and resolutions}

Three build iterations. Two fixes were routine (an \texttt{omega} for a
doubled natural equation; replacing a \texttt{field\_simp} that mangled
an \texttt{rpow} into a subtype projection by deterministic
cancellation of positive factors). The instructive one: dropping the
\emph{unused} hypothesis $0 < \alpha_2$ from the rigidity lemma
strengthened the statement but left $\alpha_2, \beta_2$ unconstrained at
the application site, and the elaborator unified them with the left-hand
values before the proof body could pin them, producing a vacuous goal.
The repair is to pass the right-hand data explicitly as named arguments
($\alpha_2 := \ldots$, $\beta_2 := \ldots$) when the only anchor
hypothesis mentions one side.

\section{What was Mathlib, what was new}

Mathlib: \texttt{mul\_rpow}, strict monotonicity of \texttt{rpow} in the
exponent, base injectivity \texttt{rpow\_left\_injOn}, positivity of
\texttt{Gamma}. Seabed: the closed form from the May monomial tides,
consumed here for the first time by a later tide. New: the sample-point
elimination and the packaging.

\section{Lessons}

Dropping an unused hypothesis is not free when it was the elaborator's
only anchor for an implicit argument: check the call sites. The
two-point sampling trick is a good template for rigidity statements
about power laws; it avoids both logarithms and limits.

\section{Follow-ups}

\begin{itemize}
\item The note's full one-dimensional induction: higher moments
recover the subleading coefficients $a_j$ of a general analytic
potential, order by order. The seabed's \texttt{kth\_moment\_even}
closed forms cover the base case; the inductive step needs the scaled
covariance pairing.
\item The semi-quasi-homogeneous weight recovery
($q_i = (\ell(2e_i) - \ell(0))/2$) once a multivariate moment layer
exists.
\end{itemize}

\end{document}
